\documentclass[10pt]{article}
\usepackage[margin=0.75in]{geometry}
\usepackage{amsmath,amssymb}
\usepackage{booktabs}
\usepackage{enumitem}
\setlist{nosep,leftmargin=1.4em}
\pagestyle{empty}

\newcommand{\E}{\mathbb{E}}
\newcommand{\HSIC}{\mathrm{HSIC}}
\newcommand{\argmin}{\operatorname*{arg\,min}}

\begin{document}
\begin{center}
{\large\textbf{Cross-Level Business-Metric Diagnosis: Association \& $Y\!\mid\!X$ Shift}}\\[2pt]
{\small Two complementary, kernel/causal-objective tools --- no explicit metric formulas required}
\end{center}
\vspace{-4pt}

Business metrics live at different levels (order, user, merchant) with no
explicit relationships. Two questions, two tools:

\vspace{4pt}
\noindent\textbf{(A) Which metrics move together? --- HSIC (association linkage).}
For metric pairs $(a,b)$ with RBF Gram matrices $K,L$ and centering
$H=I-\tfrac1n\mathbf 1\mathbf 1^\top$,
\[
\widehat{\HSIC}(a,b)=\tfrac{1}{(n-1)^2}\operatorname{tr}(KHLH)\ \ge 0,
\qquad =0 \iff a\perp\!\!\!\perp b .
\]
Symmetric, non-linear, assumption-light; large $\widehat{\HSIC}$ = a linkage
edge. It recovers the chain and isolates unrelated metrics (\emph{no direction
claimed}).

\vspace{3pt}
\noindent\textbf{(B) Did a target's $Y\!\mid\!X$ change, and on which $X$? ---
PO-risk (one model per target).}
With batch $W$ as ``treatment'', cross-fit $\mu(x)=\E[Y|x]$, $e(x)=\Pr(W{=}1|x)$;
pseudo-outcome $\psi=(Y-\mu)(W-e)$; $\tau=\argmin_f\E[(\psi-f(X))^2]$;
\[
\mathrm{PO\text{-}risk}=\E[\tau(X)^2]\ (>0 \Rightarrow P(Y\!\mid\!X)\text{ shifted}),
\qquad
\phi_j=\mathrm{PO\text{-}risk}(X)-\mathrm{PO\text{-}risk}(X_{-j}).
\]
One meta-learner takes all $X$ jointly; LOCO $\phi_j$ localizes the change ---
cost $O(\#\text{targets})$, not $O(p^2)$ pairwise models.

\vspace{6pt}
\begin{center}
\begin{tabular}{@{}p{0.46\textwidth}p{0.46\textwidth}@{}}
\toprule
\textbf{(A) HSIC linkage} & \textbf{(B) PO-risk} $Y{=}$GMV (test $p=0.010$)\\
\midrule
\begin{tabular}{@{}llr@{}}
$a$ & $b$ & $\widehat{\HSIC}$\\
\midrule
delivery & rating   & 0.055\\
delivery & gmv      & 0.046\\
rating   & n\_orders& 0.062\\
n\_orders& gmv      & \textbf{0.070}\\
$\ast$ w/ quantity/discount & & $\le 0.002$\\
\end{tabular}
&
\begin{tabular}{@{}lr@{}}
metric $x_j$ & LOCO $\phi_j$\\
\midrule
\textbf{rating} (driver) & \textbf{+0.106}\\
orders   & +0.046\\
delivery & +0.039\\
promo    & +0.032\\
tenure   & +0.028\\
\end{tabular}
\\
\midrule
Chain \texttt{delivery--rating--n\_orders--gmv} recovered; noise metrics isolated.
& $P(Y|X)$ shifted; change loads on \texttt{rating}.\\
\bottomrule
\end{tabular}
\end{center}

\vspace{4pt}
\noindent\textbf{Use together.} HSIC maps the metric linkage across levels; for
any target metric, PO-risk says whether its $Y\!\mid\!X$ relationship shifted and
which linked metric carries it --- so a change in a merchant-level metric can be
localized to an order-level metric, without modelling every pair and without a
causal-proportion claim.

\vspace{3pt}
\noindent\textbf{Central point.} This feature selection for distribution shift
(FSDS) procedure extends beyond two batches, unlocking multi-layer feature
selection and localization in both feature stores and concurrent metric systems.
In many circumstances the vague associations between metrics are not yielded by
business resource-intensive models (which require a lot of man-days from the
business requirement to deployment) and are not deployed in production systems
--- this is where another value of our FSDS procedure lies.

\end{document}
